\section{Related Work}
\label{sec:related-work}

\textbf{Contamination evidence.} Sainz~\emph{et al.}~\cite{sainz2023contamination} argue contamination measurement should be standard practice, not assumed away; Balloccu~\emph{et al.}~\cite{balloccu2024leakcheat} support this empirically, tracing leaked benchmark references into closed-source model cards. LLMs reproduce training sequences at rates that grow with model scale and duplication~\cite{carlini2023quantifying}; Dong~\emph{et al.}~\cite{dong2026generalization} extend this to code generation. These studies establish contamination is real and detectable after training happens; none offers a protocol that prevents it at measurement time, the gap this work addresses for test generation.

\textbf{Dynamic evaluation.} A growing line of work responds to static-benchmark contamination by generating evaluation instances at test time. DyVal~\cite{zhu2024dyval} constructs reasoning problems from dynamic graphs; LiveBench~\cite{white2025livebench} refreshes monthly against updated sources; BeyondBench~\cite{srivastava2026beyondbench} generates algorithmic problems with correctness guarantees; MATHEMAGIC~\cite{obrien2025mathemagic} applies the same principle to mathematics; LLM-as-an-Interviewer~\cite{kim2025interviewer} replaces static sets with dialogue-driven evaluation. Each resists the contamination failure mode static benchmarks exhibit; none targets unit-test generation, where a generated test is assessed by execution rather than comparison to a reference string.

\textbf{Dynamic and contamination-aware benchmarking for code.} Closer to the present setting, DyCodeEval~\cite{chen2025dycodeeval} and the benchmark-usefulness study of Guan~\emph{et al.}~\cite{guan2025benchmarkuseful} apply dynamic regeneration to code-reasoning benchmarks; Code2Bench~\cite{zhang2025code2bench} scales dynamic construction directly from source repositories; LiveCodeBench~\cite{jain2025livecodebench} draws problems from programming contests held after training cutoffs. Automated unit-test generation predates LLMs: EvoSuite~\cite{fraser2011evosuite} and Pynguin~\cite{lukasczyk2022pynguin} synthesise suites by evolutionary search, and recent studies evaluate LLMs on the task directly~\cite{schafer2024empirical, siddiq2024junit}. TestEval~\cite{wang2024testeval} and Haroon~\emph{et al.}~\cite{haroon2026evolution} evaluate LLM test generation, the latter tracking quality as the codebase evolves; MuCoCo~\cite{chou2026mucoco} evaluates code-LLM consistency via mutation-style perturbation, a close methodological parallel here. The static comparator adopted in this work, UnLeakedTestBench (ULT)~\cite{huang2026ult}, decontaminates real-world functions and constructs PLT, a deliberately leaked twin, so any contamination-consistent score gap becomes observable, a design this work reuses. Codex/HumanEval~\cite{chen2021codex} established code-generation benchmarking but is reported contaminated in later models; BigCodeBench~\cite{zhuo2025bigcodebench} serves here as an external code-generation proxy for construct validity.

\textbf{The gap.} No prior dynamic-benchmarking study combines template-based, seeded regeneration with a mutation-testing oracle for unit-test generation specifically, and none pairs an observational contamination comparison (dynamic versus a decontaminated static benchmark and its leaked twin) with a causal one (direct fine-tuning on the static benchmark's own ground truth, isolating contamination-driven inflation from generic capability gain). This work supplies both, together with a construct-validity check against an external, non-test-generation proxy, to establish not only that DyMuT resists contamination but that it continues to measure a capability that an independent evaluation recognises. The supporting report's own literature review tabulates this comparison in full, against nine dynamic-benchmarking systems and a second group of LLM test-generation benchmarks, each with what it cannot do that motivated a design choice here.
